%% conformal-antenna-diffopt.5: rewrite around the physical 2.4 GHz bent patch on a fixed
%% 1.6 mm FR-4 board (copper-only design), its mesh-refinement study and the frequency
%% space-mapping correction. See changelog.md (critic notes and operator inputs mapped to
%% changes) and provenance.md (every number to its committed source).
%% v1-v4 and their critic siblings are unchanged.

\documentclass{anvil-paper}
\usepackage{xurl}
\usepackage{array}

\expandafter\def\expandafter\UrlBreaks\expandafter{\UrlBreaks\do\-\do\_}
\Urlmuskip=0mu plus 1mu
\renewcommand{\bibfont}{\small}
\setlength{\bibsep}{2pt plus 1pt}

\newcommand{\anvilfig}[3][0.98]{%
  \IfFileExists{figures/#2}%
    {\includegraphics[width=#1\linewidth]{figures/#2}}%
    {\fbox{\parbox{#1\linewidth}{\centering\vspace{5em}%
      \texttt{figures/#2}\\[0.5em]%
      {\small placeholder: render via \texttt{figures/src/#3}}%
      \vspace{5em}}}}}
\newcommand{\PL}{P2/L1}

\title{Shape-Adjoint Design of a Curved Conformal 2.4\,GHz Patch Antenna
on a Fixed 1.6\,mm Board with a Differentiable Open-Radiator FEM}
\author{Robb Walters\\
  {\normalsize Independent Researcher}}
% NOTE: Crutcher Dunnavant to be offered coauthorship at review (operator, 2026-10-13)
\date{\today}

\begin{document}
\maketitle

\begin{abstract}
We design a probe-fed 2.4\,GHz patch antenna wrapped on a cylinder of
radius $100$\,mm, keeping the board a standard uniform $1.6$\,mm FR-4
shell and changing only the etched copper. The design tool is a
differentiable three-dimensional N\'ed\'elec finite-element method (FEM)
on unstructured tetrahedral meshes, whose discrete shape adjoint passes
through the full open-radiator forward model (box UPML absorber, lossy
substrate and a lumped feed port whose position and width move), at one
LU factorization per frequency. Thirty-eight smooth copper degrees of
freedom (outline splines, corner bumps, feed inset and width) are applied
by an exact on-cylinder morph, so the board stays $1.6$\,mm thick. A
mesh-refinement study shows that the antenna's resonance is not converged
on any affordable mesh: it rises under every $h$- and $p$-refinement, and
a design optimized on the L1 mesh resonates about $45$--$48$\,MHz higher
with second-order elements. That detune is close to a pure frequency
translation, so we correct it by frequency space mapping and verify on
the finest computed model, second-order elements on the second mesh level
(\PL). There the design reaches a worst in-band $|S_{11}|$ of
$-10.78$\,dB and a $94.9$\,MHz $-10$\,dB band ($2.3938$--$2.4887$\,GHz)
that covers the $2.400$--$2.4835$\,GHz ISM band. On the finest
first-order mesh it reaches $-8.80$\,dB. Richardson estimates for the initial antenna
place its converged resonance a further $12.7$--$48.4$\,MHz up; the design
has tracked the initial antenna's shifts to within $3$\,MHz, and every
estimate exceeds the $6.2$\,MHz margin at the lower band edge, so the
converged design probably misses at $2.400$\,GHz. The feed is an idealized $4.3$\,mm port strip.
\end{abstract}

\section{Introduction}
\label{sec:intro}

Gradient-based inverse design has become routine in nanophotonics, where
differentiable solvers such as \textsc{ceviche} \citep{hughes2019forward},
the Meep adjoint solver \citep{hammond2022meepadjoint}, SPINS-B
\citep{su_spinsb} and FDTDX \citep{mahlau2026fdtdx} optimize a permittivity
density on a structured grid. The same density approach has been carried
to metallic antennas: \citet{hassan2014topology} optimize the conductivity
of a microstrip patch on a finite-difference time-domain (FDTD) grid with an
adjoint gradient and more than $20{,}000$ design variables. Shape
derivatives of scattering parameters for finite-element models with ports
are older still \citep{webb2000design,webb2002design,nikolova2004adjoint}
and are available in commercial tools
\citep{ansys_hfss_derivatives,comsol_sensitivity_method}.

This paper uses a differentiable finite-element solver, geode-fem, to
design a specific antenna under manufacturing constraints and asks how far
the result can be trusted. The antenna is a 2.4\,GHz probe-fed patch bent
around a cylinder of radius $R = 100$\,mm. The board is a uniform
$1.6$\,mm FR-4 shell and only the copper changes, so a design is a copper
outline and a feed position that a board house could etch on standard
stock. The body-fitted tetrahedral mesh follows the curved conductor
exactly, and the design variables move it.

Two findings shape the paper. First, the forward model is not converged on
any mesh we can afford: the resonance of the initial antenna rises by
$85$\,MHz from the coarsest to the finest computed model, and the order of
convergence is falling along the ladder, so the converged resonance is
bounded below but not above (Section~\ref{sec:convergence}). A design
optimized on a coarse mesh therefore lands off band on a finer one.
Second, that discrepancy is close to a pure frequency translation, which a
frequency-only form of space mapping \citep{bandler1994space,bandler2004space}
corrects with one or two refits (Section~\ref{sec:procedure}).

The contributions are:
\begin{itemize}
  \item A discrete shape adjoint for a radiating antenna model on
  unstructured tetrahedra: N\'ed\'elec elements, a box UPML, a lossy
  substrate and a lumped port whose position and width are design
  variables, with element Jacobians from forward-mode automatic
  differentiation and one LU factorization per frequency shared by the
  forward and adjoint solves (Section~\ref{sec:method}). The gradient
  matches central finite differences of the full pipeline to a relative
  $2.2\times10^{-5}$ or better on every run.
  \item A copper-only design parametrization on a fixed curved board:
  $38$ smooth degrees of freedom applied by an exact morph that keeps every
  substrate node on its own cylinder (Section~\ref{sec:freedom}).
  \item A design verified against refinement: worst in-band $-10.78$\,dB
  and a $94.9$\,MHz $-10$\,dB band on the finest computed model, with the
  finer first-order mesh, an independent sibling mesh and an extrapolation
  of the converged response reported alongside, including the likely miss
  at the lower band edge (Section~\ref{sec:results}).
  \item Two ablations: without the frequency correction the same procedure
  gives $-4.76$\,dB on the finest model, and an idealized design that also
  varies the board thickness reaches $-9.16$\,dB there
  (Section~\ref{sec:ablations}).
\end{itemize}

\section{Related Work}
\label{sec:related}

\paragraph{Finite-element S-parameter sensitivities with ports.}
Design sensitivities of the scattering matrix for tetrahedral edge
elements were derived by \citet{webb2000design}, extended to the frequency
response by \citet{webb2002design} and used in gradient optimization of
microwave devices by \citet{nair2003optimization}. Adjoint and
self-adjoint formulations for S-parameter sensitivities are reviewed by
\citet{nikolova2004adjoint} and \citet{nikolova2006network}, and continuum
shape sensitivity for high-frequency devices is given by
\citet{choi2011generalized}. HFSS computes analytic S-parameter derivatives
with respect to geometric parameters by reusing the LU factors
\citep{ansys_hfss_derivatives,vardapetyan2008sensitivity}, and COMSOL
documents an adjoint sensitivity method and shape optimization of a
waveguide filter and a diplexer on a deformed mesh
\citep{comsol_sensitivity_method,jensen2022comsolshape}. These examples are
closed waveguide devices with a few geometric parameters. For antennas,
\citet{ghassemi2013antenna} evolve a microstrip antenna with an FEM adjoint
sensitivity and a small set of control vertices. Our model differs in
being an open radiator (box UPML, lumped port, lossy FR-4) with $38$
smooth boundary degrees of freedom and a moving feed.

\paragraph{Automated shape differentiation.}
Shape derivatives of finite-element models can be generated
automatically: in UFL and Firedrake \citep{ham2019automated,ham2020,paganini2021fireshape}
and in NGSolve \citep{gangl2021fully}. \citet{wang2011} give a discrete
Maxwell shape adjoint on unstructured meshes in two dimensions, in the time
domain and without ports. geode-fem obtains the element Jacobians by
forward-mode differentiation of its closed-form element kernels inside a
tensor stack \citep{burn}, rather than by symbolic shape calculus.

\paragraph{Structured-grid design of metallic antennas.}
Density-based topology optimization of conductors on a fixed grid goes
back at least to \citet{erentok2011topology}. \citet{hassan2014topology}
and \citet{hassan2014patch} optimize the patch and ground plane of
coax-fed microstrip antennas with an FDTD adjoint, and
\citet{nikolova2004sensitivity,nikolova2006scattering} give FDTD
sensitivities with respect to shape parameters. A plain Yee grid
staircases a curved conductor, which produces $O(1)$ field errors next to
the boundary \citep{cangellaris1991analysis,haggblad2014accuracy}, but
structured grids need not staircase: locally conformal FDTD has handled
curved PEC for decades \citep{dey1997locally,mittra2000conformal},
including uniformly stable second-order schemes
\citep{zagorodnov2003uniformly,benkler2006new}. In differentiable
photonics, subpixel-smoothed projection gives the Meep density adjoint
subpixel-accurate binarized boundaries \citep{hammond2025unifying,romano2026ssp2},
though the material-grid adjoint excludes metallic media
\citep{meepadjointdocs}; Tidy3D offers autograd shape gradients and
conformal PEC averaging in its forward solver \citep{tidy3dautograd}; and
\citet{hooten2025} and \citet{liu2024imagerep} map shape parameters to a
fixed grid. Precomputed numerical Green's functions make pixelated antenna
design fast without gradients \citep{sun2025pngf}. We did not find a
conformal FDTD adjoint for the shape of a curved conductor, and we do not
compare against any of these methods here (Section~\ref{sec:limits}).

\paragraph{Integral-equation shape sensitivity.}
Boundary- and integral-equation methods carry recent PEC shape-sensitivity
work: time-domain boundary-integral shape derivatives
\citep{takahashi2025tdbie}, freeform tubular scatterers
\citep{arens2026tubular} and isogeometric coaxial baluns
\citep{balouchev2024balun}.

\paragraph{Space mapping.}
Space mapping aligns a cheap coarse model with an accurate fine model and
optimizes the coarse one \citep{bandler1994space,bandler2004space}. We use
its simplest form, a frequency shift fitted between two discretizations of
the same solver.

\paragraph{Position of this work.}
The closest prior work is the FEM S-parameter sensitivity line above,
which already reuses factorizations and differentiates with respect to
geometry. What we add is the combination, for an open radiator, of a box
UPML, a lumped port whose position and width move, a lossy substrate and a
smooth many-parameter copper basis on a curved board, together with a
design verified against mesh and order refinement. We do not claim a first
EM shape adjoint, and we make no performance comparison with structured-grid
methods.

\section{Method}
\label{sec:method}

\begin{figure}[t]
  \centering
  \anvilfig{fig1_geometry.pdf}{fig1_geometry.py}
  \caption{The fixture and the final copper design. (a) Section at $y = 0$,
  to scale: the FR-4 board (green) on the $R = 100$\,mm cylinder with the
  patch on top, inside the air box and the $25$\,mm box UPML (hatched).
  The bend is blended to zero at the air-box wall, so the UPML stays a
  rectangular box. (b) Section through the feed at $s = 0$, with the
  vertical axis stretched four times: patch, $1.6$\,mm FR-4, ground, and
  the lumped-port strip across the gap, final (filled) and initial
  (dashed). (c) Top view in the chart $(s, y)$: the initial $38\times29$\,mm
  patch (dashed), the final copper outline (filled), and the initial and
  final port strips. The board outline and thickness are the same in both.
  Drawn from \texttt{baseline\_2g4\_L1.toml} and
  \texttt{geometry/rbD\_L1\_copper\_38\_sm45.json}.}
  \label{fig:geometry}
\end{figure}

\subsection{The open-radiator forward model}
\label{sec:pencil}

The forward problem is the driven frequency-domain curl-curl equation,
discretized with N\'ed\'elec edge elements of order $p = 1$ or $p = 2$ on
an unstructured tetrahedral mesh. Lengths are in millimetres and the
solver works with $\omega = k_0 = 2\pi f/c$ in rad/mm. The assembled system
is the complex-symmetric pencil
\begin{equation}
  A(\omega)\,\mathbf{x}
  = \Big[\,K(\nu) - \omega^2 M(\varepsilon) + \tfrac{j\omega}{Z_s}\,S_p\,\Big]\,\mathbf{x}
  = \mathbf{b},
  \label{eq:pencil}
\end{equation}
with curl-curl stiffness $K$, mass matrix $M$ and lumped-port surface
term $S_p$. The radiating boundary is a box UPML: per-element tensor
weights $\Lambda$ enter both $K$ (as $\nu = \Lambda^{-1}$) and $M$ (as
$\varepsilon = \varepsilon_r \Lambda$), with stretching
$s = 1 - j\sigma_0 (d/w)^2/\omega$ per axis. The FR-4 has
$\varepsilon_r = 4.4$ and $\tan\delta = 0.02$, and the conductors are
perfect electric conductors (PEC). The port is a rectangle of width $w$
along $y$ and length $l = 1.6$\,mm across the substrate gap, with sheet
impedance $Z_s = R\,w/l$ and $R = 50\,\Omega$ (entered in the solver as
$50/\eta_0$). The reflection coefficient follows from the port voltage $V$:
\begin{equation}
  S_{11} = \frac{Z - R}{Z + R},\qquad Z = \frac{V}{I},\qquad
  I = \frac{2V_{\mathrm{inc}} - V}{R}.
  \label{eq:s11}
\end{equation}
Every evaluation checks passivity, $|S_{11}| \le 1$.

\subsection{Shape adjoint with one factorization per frequency}
\label{sec:adjoint}

The design moves mesh node coordinates $X$, so the gradient needs
$\partial K/\partial X$, $\partial M/\partial X$, $\partial S_p/\partial X$
and $\partial \mathbf{b}/\partial X$. geode-fem evaluates the closed-form
element kernels in first-order dual arithmetic: seeding the tangent of one
node coordinate returns the exact partial derivatives of every local
matrix entry with respect to it. The UPML weights $\Lambda$ are held fixed
under node motion, so the tangent flows through geometry only. Because
$A(\omega)$ is complex symmetric, the adjoint system uses the same
operator, and one LU factorization per frequency serves both the forward
and the adjoint solve. The nodal gradient at each frequency is
$\lambda^{\top}(\partial A/\partial X)\,\mathbf{x}$ plus the explicit
dependence of the port drive and of the port-voltage functional on the
geometry.

When a design variable moves the feed, the port faces move: $w$, $l$, the
drive and the voltage functional all depend on $X$. The moving-port terms
are differentiated together with the volume terms, and the port is rebuilt
on every moved mesh. A unit test of the moving-feed gradient on a curved
mesh, with a current source that varies across the port so that the
explicit term is not negligible, matches central finite differences to
below $10^{-8}$; without the explicit term it misses by a relative error
of about $3$.

\subsection{Copper design freedom on a fixed curved board}
\label{sec:freedom}

The board is not a design variable. Each substrate node is described in a
chart $(s, y)$, with $s = R\,\mathrm{atan2}(x, z - z_{\mathrm{axis}})$ the
mid-plane arc length and $y$ the bend axis. A design applies a chart
displacement $(\Delta s, \Delta y)$ and moves the node to angle
$\phi_0 + \Delta s/R$ and height $y_0 + \Delta y$ at its own base radius.
The radial column through the slab therefore rotates together, the FR-4
stays a uniform $1.6$\,mm shell, the patch stays on the $R + h/2$ surface
and the ground on $R - h/2$. Every run records the largest change of any
substrate node's radius ($\le 1.4\times10^{-14}$\,mm) and the slab thickness
($1.600$\,mm everywhere). Air nodes follow a harmonic (P1-Laplace)
extension of the motion; the UPML shell and the outer boundary are fixed.
The design gradient chains the nodal adjoint gradient through the exact
Jacobian of this map.

Table~\ref{tab:dofs} lists the $38$ degrees of freedom. They are smooth
functions of the chart alone, so the same coefficient vector defines the
same copper on every mesh of the family, which is what lets a design be
carried between discretizations. The basis cannot open a hole in the
patch: slots (U-slot, E-patch) change the topology of the conductor and are
outside this parametrization.

\begin{table}[t]
  \centering
  \caption{The $38$ copper degrees of freedom. Units are per unit of the
  coefficient. The board thickness has no degree of freedom.}
  \label{tab:dofs}
  \small
  \begin{tabular}{@{}l c >{\raggedright\arraybackslash}p{0.66\linewidth}@{}}
    \toprule
    Group & Count & What moves \\
    \midrule
    Outline & $32$ & $8$ clamped cubic B-splines per patch edge ($\pm y$, $\pm s$); each moves its edge outward in the surface by $1$\,mm, ramped in from the feed band and decaying over the $12$\,mm substrate margin \\
    Corners & $4$ & one bump per corner, moving it outward along the bisector by $1$\,mm (negative values chamfer), support $1 - \mathrm{smoothstep}(d/6\,\mathrm{mm})$ \\
    Feed & $2$ & the port's position along $y$ (its inset from the $+y$ edge) and its width $w$ along $y$, each $1$\,mm \\
    \bottomrule
  \end{tabular}
\end{table}

\subsection{Objective, optimizer and guards}
\label{sec:objective}

The objective is a soft maximum of $|S_{11}|^2$ over $n_f = 5$ band
frequencies,
\begin{equation}
  G(\mathbf{c}) = \frac{1}{\beta}\ln \sum_{f} \exp\!\big(\beta\,|S_{11}(f;\mathbf{c})|^2\big),
  \qquad \beta = 25,
  \label{eq:objective}
\end{equation}
over the coefficient vector $\mathbf{c}$. It is minimized by L-BFGS
(memory $8$) with Armijo backtracking and a cap of $1$ coefficient unit
per step. A per-element guard rejects any trial step that would shrink a
tetrahedron's signed volume below $0.25$ of its base value. In the runs
reported here the guard, not the gradient, ends the descent: every run
stops at the volume budget (or on a plateau) with the gradient norm still
nonzero, so the designs are guard-limited rather than stationary.

Before and after each run the full pipeline gradient is compared with a
central finite difference along the steepest-descent direction and along
random directions in all $38$ coefficients. The worst relative error over
all runs is $2.2\times10^{-5}$ (tolerance $5\times10^{-3}$); a per-group
check of the outline, corner, inset and width coefficients on the two
coarsest meshes passes to $1.9\times10^{-7}$. The adjoint is implemented
for $p = 1$ only, so optimization runs at $p = 1$ and second-order
elements enter through verification.

\section{Fixture and Discretization}
\label{sec:fixture}


\paragraph{Antenna.}
The initial antenna is a $38\times29$\,mm PEC patch (width $38$\,mm along
$s$, resonant length $29$\,mm along $y$) on $1.6$\,mm FR-4 with a $12$\,mm
substrate and ground margin, fed by a lumped-port strip $1.6$\,mm wide
inset $7.0$\,mm from the $+y$ edge (Figure~\ref{fig:geometry}). The flat
mesh is wrapped onto a cylinder of radius $100$\,mm about an axis
parallel to $y$; the $62$\,mm substrate then subtends $\pm0.31$\,rad, with a
sagitta of $1.8$\,mm at the patch edge. A $40$\,mm air margin separates the
board from a $25$\,mm box UPML. The target is the $2.400$--$2.4835$\,GHz ISM
band at $-10$\,dB.

\paragraph{Discretization ladder.}
Four graded meshes L0--L3 carry the same geometry with $38{,}650$,
$150{,}561$, $370{,}780$ and $844{,}238$ edges, and second-order elements
run on L0 and L1 ($204{,}824$ and $804{,}460$ unknowns). We write \PL{} for
$p = 2$ on L1. The UPML conductivity $\sigma_0 = 1$ was chosen from a
ladder of $\sigma_0$ and shell thickness on L1, where the residual
sensitivity of $S_{11}$ to the absorber setting is $|\Delta S_{11}| \le 0.021$.

\paragraph{Bending.}
At this radius the bend barely moves the resonance: the bent dip lies
$1.0$ to $1.6$\,MHz below the flat dip on every level. The curvature
matters for representing the conductor and for the design freedom, not as
a large detune.

\subsection{The forward model is not converged}
\label{sec:convergence}

\begin{figure}[t]
  \centering
  \anvilfig{fig2_convergence.pdf}{fig2_convergence.py}
  \caption{Resonance (the $|S_{11}|$ dip) of the initial bent patch under
  refinement, against the number of unknowns. Left: $p = 1$ on L0--L3 and
  $p = 2$ on L0 and L1. The dip rises under every refinement step; the
  dotted line is the finest computed value, $2.3566$\,GHz at \PL, a lower
  bound on the converged dip under that monotone trend. Right: Richardson
  extrapolations of the $p = 1$ ladder and of the $p = 2$ pair at the
  three orders shown ($1$, the local order of the finest triple, and a
  least-squares slope). They are estimates under assumed orders, not
  bounds. The shaded band is the ISM band. Data:
  \texttt{convergence.toml}.}
  \label{fig:convergence}
\end{figure}

Figure~\ref{fig:convergence} shows where the initial antenna resonates on
each model. The dip rises under every $h$- and $p$-refinement step, from
$2.2717$\,GHz on L0 to $2.3422$\,GHz on L3 and $2.3566$\,GHz at \PL. At a
similar number of unknowns \PL{} lies $14.5$\,MHz above L3, so the
second-order model is the more accurate one, and the finest computed dip
bounds the converged dip from below. The upper end is not determined. The
local order of the dip frequency falls along the $p = 1$ ladder, from
$1.18$ on L0--L1--L2 to $0.50$ on L1--L2--L3, so the levels are not in the
asymptotic range. Extrapolating the $p = 2$ pair gives $2.3693$,
$2.3779$ and $2.4050$\,GHz at orders $1.49$, $1$ and $0.50$, that is
$12.7$, $21.3$ and $48.4$\,MHz above \PL. Order $1$ is the rate expected
from the $r^{-1/2}$ field singularity at the patch edge, but we have no
measurement that selects it.

For design this has two consequences. A design centred on any computed
model will resonate higher on a finer one. And since the in-band worst
case of the initial antenna improves with fidelity but is still $-2.42$\,dB
at \PL, the optimization problem is real at every fidelity we can compute.

Costs set the ladder. On a $28$-core workstation an L1 forward solve takes
about $44$\,s and a band gradient at three frequencies $67$\,s, at
$6$--$8$\,GiB peak memory. A \PL{} solve takes $753$\,s and $132$\,GiB on
a $32$-vCPU cloud host, and an L3 solve about $8$\,min and $68$\,GiB. We
therefore optimize on L1 or L2 and verify on L3 and \PL.

\section{Design Procedure}
\label{sec:procedure}

\paragraph{Coarse optimization.}
A cold start on L1 optimizes the band directly and reaches a worst-of-band
$-10.04$\,dB on L1 after $24$ steps, ending when the volume guard binds in
a tetrahedron between the feed and the $+y$ edge. With \PL{} the same
copper gives only $-4.76$\,dB (Figure~\ref{fig:response}b): the design
resonates higher on the finer model.

\paragraph{Frequency space mapping.}
We model the fine response as the coarse response shifted in frequency,
$S_{\mathrm{fine}}(f) \approx S_{\mathrm{coarse}}(f - \delta)$, and fit
$\delta$ by least squares over the fine model's sampled points, with the
coarse response interpolated on a dense coarse sweep. For the cold L1
optimum $\delta = 47.85$\,MHz, and the root-mean-square complex misfit
falls from $0.23$--$0.38$ unshifted to $0.018$--$0.025$ shifted, over the
copper designs fitted. The shift is within $3$\,MHz of the initial patch's
own L1-to-\PL{} shift ($44.8$\,MHz), so the copper motion adds almost no
detune of its own. The corrective step re-optimizes on the coarse model
for the band shifted down by $\delta$ and checks the result at \PL. One
refit was needed: the first corrected design gave $-9.96$\,dB at \PL, the
refitted shift was $44.94$\,MHz, and the second corrected design is the
final one.

\paragraph{Feed re-basing.}
Moving the feed toward the $+y$ edge compresses the $5$--$7$\,mm of mesh
between them, and the cold runs on every level stop there at the volume
guard. We then regenerate the mesh with the port at its moved position and
width and carry the design over, shifting the two feed coefficients by the
change of the fixture's port. From an original mesh this carries the
copper exactly. Between two regenerated meshes it does not: the feed
coefficients' support also reaches the $+y$ edge there, and the re-base
drops that motion, so the target mesh carries a slightly different copper.
We report that difference wherever it applies (Section~\ref{sec:results}).
After re-basing, the optimizer widens the patch from $38$ to about
$49.8$\,mm and the port from $1.6$ to about $4.3$\,mm, until the guard binds
in the substrate margin $5$--$6$\,mm outside a radiating edge.

Table~\ref{tab:chain} lists the chain that produced the final design.

\begin{table}[t]
  \centering
  \caption{The optimization chain of the final design. ``Worst'' is the
  worst of the band points on the band each run optimized (shifted down
  by $\delta$ where stated), on the mesh it ran on. All runs are at
  $p = 1$ on L1 meshes; ``re-based'' means the mesh was regenerated at the
  moved port. The wall times are on cloud hosts of $16$ or $32$ vCPU.}
  \label{tab:chain}
  \small
  \begin{tabular}{@{}l l c c c l c@{}}
    \toprule
    Step & Start & $\delta$ (MHz) & Worst (dB) & Steps & Terminal & Wall \\
    \midrule
    1 cold & initial patch & $0$ & $-1.47 \to -10.04$ & $24$ & guard (feed) & $0.8$\,h \\
    2 re-based & step 1 & $0$ & $-9.97 \to -10.03$ & $4$ & plateau & $0.5$\,h \\
    3 corrected & step 2 & $47.85$ & $-4.75 \to -10.25$ & $26$ & guard (widening) & $0.8$\,h \\
    4 refit, re-based & step 3 & $44.94$ & $-10.26 \to -10.78$ & $9$ & guard (widening) & $0.7$\,h \\
    \bottomrule
  \end{tabular}
\end{table}

\section{Results}
\label{sec:results}

\begin{figure}[t]
  \centering
  \anvilfig{fig3_response.pdf}{fig3_response.py}
  \caption{(a) Return loss of the final design at \PL{} (blue, $41$
  points) and on L3 (teal, dashed, $41$ points), and of the initial patch
  at \PL{} (gray, sampled near its dip and at nine in-band points). The
  dashed orange line is $-10$\,dB and the shaded band is the ISM band. The
  arrow starts at the lower $-10$\,dB edge at \PL{} and has the length of
  the largest extrapolated further upshift ($48.4$\,MHz); the smallest is
  $12.7$\,MHz (Section~\ref{sec:convergence}). (b) The cold L1 optimum
  (step~1 of Table~\ref{tab:chain}) on L1, where it was optimized, its
  nine in-band points at \PL, and the L1 curve shifted by the fitted
  $47.85$\,MHz. Data: \texttt{copper/verify/} and \texttt{convergence/}.}
  \label{fig:response}
\end{figure}

\paragraph{Finest computed model.}
At \PL{} the final design has a worst $|S_{11}|$ of $-10.78$\,dB over nine
points spanning the ISM band, and the $41$-point sweep gives a contiguous
$-10$\,dB band from $2.3938$ to $2.4887$\,GHz, $94.9$\,MHz wide, which covers
the whole $83.5$\,MHz ISM band (Figure~\ref{fig:response}a,
Table~\ref{tab:verify}). The band is centred to $0.5$\,MHz. The margins are
small: the $-10$\,dB band starts $6.2$\,MHz below $2.400$\,GHz and ends
$5.2$\,MHz above $2.4835$\,GHz. This verification runs on the mesh the
design was optimized on, so the verified copper is the optimized copper
(they agree to about $10^{-3}$\,mm), and its minimum volume ratio is the
guard value $0.250$. As built, the patch is $28.80$\,mm long and $49.75$\,mm
wide at the feed, the port strip is $4.27$\,mm wide and $5.07$\,mm from the
$+y$ edge, and the board is $1.600$\,mm thick everywhere.

\paragraph{Finest first-order mesh.}
On L3 the same copper gives $-8.80$\,dB and covers $89\%$ of the band. Its
$-10$\,dB band is about as wide ($94.6$\,MHz) but sits about $14$\,MHz lower,
close to the initial patch's $14.5$\,MHz L3-to-\PL{} offset. With $95$\,MHz
of bandwidth against an $83.5$\,MHz band there are only about $11$\,MHz of
slack, so a design centred for one model is off centre on the other. The
L3 mesh of the widened patch has thin air elements next to the slab
margin, which the default air morph compresses to a volume ratio of
$0.10$; moving the air on cylinders instead raises it to $0.24$ and changes
$|S_{11}|$ by $0.01$\,dB, so the L3 shortfall is not a mesh-quality
artifact.

\paragraph{Where the converged design probably falls.}
\PL{} is the better reference: the dip rises under refinement and \PL{}
lies above L3 at a similar size, so L3 misses on the less accurate side.
The same data put the converged response above \PL. The extrapolated
further rise of the initial patch's dip is $12.7$--$48.4$\,MHz
(Section~\ref{sec:convergence}), and the copper design's coarse-to-\PL{}
shift tracks the initial patch's to within $3$\,MHz, so a similar further
rise is expected. Every extrapolated rise exceeds the $6.2$\,MHz margin at
the lower band edge. The design meets the $-10$\,dB target across the ISM
band on the finest computed model, and on the converged estimate it
probably misses at $2.400$\,GHz.

\paragraph{Sibling mesh.}
An independently regenerated L1 mesh at the same port gives $-10.35$\,dB at
\PL, about $3$\,MHz higher. This is not a pure remeshing test: the re-base
onto that mesh also shortened the patch by $0.09$\,mm at the feed
(Section~\ref{sec:procedure}), and a shorter patch resonates higher, so the
$3$\,MHz mixes remeshing with a small design change.

\paragraph{A related design from an L2 chain.}
A second chain optimized on L2 (shift $29.33$\,MHz) reaches a design of the
same family: patch about $28.9\times49.0$\,mm, port $4.5$\,mm. Its
verifications ran on regenerated meshes whose re-base shortened the $+y$
edge by $0.47$\,mm at the feed, so they verify a related design, not the L2
optimum: $-10.10$\,dB and a $94.3$\,MHz band covering the ISM band at \PL,
with $0.7$\,MHz of margin at $2.400$\,GHz, and $-9.25$\,dB on L3. No copper
design has been verified on L3 and at \PL{} exactly as optimized on L2.

\begin{table}[t]
  \centering
  \caption{Fresh forward verifications. ``Worst'' is over the nine in-band
  points; the band columns come from dense sweeps ($41$ points,
  $2.30$--$2.60$\,GHz; $21$ points, $2.35$--$2.55$\,GHz for the related
  design). Coverage is the fraction of the ISM band below $-10$\,dB.
  ``Vol'' is the smallest tetrahedron volume ratio on the verifying mesh.}
  \label{tab:verify}
  \small
  \begin{tabular}{@{}l l c c c c@{}}
    \toprule
    Design & Model & Worst (dB) & $-10$\,dB band (GHz) & Coverage & Vol \\
    \midrule
    Initial patch & \PL & $-2.42$ & none in band & $0\%$ & $1.00$ \\
    Final & \PL & $-10.78$ & $2.3938$--$2.4887$ & $100\%$ & $0.250$ \\
    Final & L3 & $-8.80$ & $2.3801$--$2.4746$ & $89\%$ & $0.241$ \\
    Final, sibling mesh ($-0.09$\,mm) & \PL & $-10.35$ & & & $0.234$ \\
    Related L2-chain design ($-0.47$\,mm) & \PL & $-10.10$ & $2.3993$--$2.4936$ & $100\%$ & $0.299$ \\
    Related L2-chain design ($-0.47$\,mm) & L3 & $-9.25$ & & & $0.228$ \\
    \bottomrule
  \end{tabular}
\end{table}

\paragraph{Bandwidth in context.}
The closed-form formulas of \citet{jackson1991simple}, in the form given
in Jackson's course notes \citep{jackson_ece6345}, estimate the bandwidth of
a flat rectangular patch on this substrate from its space-wave,
surface-wave and dielectric $Q$ (conductor loss omitted, since the FEM model
is PEC). At a voltage standing-wave ratio of $2$ they give about $63$\,MHz
for the initial $38\times29$\,mm patch and $70$\,MHz for the
$49.8\times28.8$\,mm final patch, with a surface-wave loss of about
$10\%$ (\path{jackson_bandwidth.py} in this paper's directory). The
FEM bandwidths are larger: $84$\,MHz for the $38$\,mm cold optimum on its
own L1 mesh, and $94.9$\,MHz for the final design at \PL. The closed form
ignores the feed and the curvature, and we have not resolved the
difference. Within the FEM results, the increase from $84$ to $95$\,MHz
comes with the two standard bandwidth measures for a single patch, a
wider patch (lower radiation $Q$) and a wider feed (lower feed
inductance), rather than any added structure.

\paragraph{Cost.}
The final chain took $2.8$\,h of optimization on cloud hosts
(Table~\ref{tab:chain}). The last run made $125$ driven solves in
$40$\,min; its first gradient (five frequencies, forward and adjoint on one
LU each) took $104$\,s at $6.2$\,GiB. The $41$-point \PL{} sweep of the
final design took $5.4$ to $7.7$\,min per frequency at $151$\,GiB peak
memory on a $32$-vCPU host, in three shards.

\section{Ablations}
\label{sec:ablations}

\paragraph{Without the frequency correction.}
The cold L1 optimum (step~1) is the same procedure without space mapping.
It reaches $-10.04$\,dB on L1 and $-4.76$\,dB at \PL. With the fitted shift
and one refit, the final design reaches $-10.78$\,dB at \PL. On this
comparison the correction, not additional design freedom, accounts for
most of the $6$\,dB. A second coarse model gives the same picture with a
smaller shift: a design optimized on L2 with a $30$\,MHz hedge gives
$-9.92$\,dB at \PL, and the fitted L2 shift is $29.33$\,MHz.

\paragraph{Varying the board thickness.}
In an idealized ablation the design freedom also moves the patch radially,
so the FR-4 under the patch thickens from $1.6$\,mm to about $3.7$\,mm on
average (up to $4.7$\,mm), with the feed left at the bottom of a $1.6$\,mm
crater. Such a board cannot be made from standard stock. Optimized on L2
without a frequency correction, the best such design reaches $-14.07$\,dB
on L2, $-12.47$\,dB on L3 and $-9.16$\,dB at \PL, missing the target by
$0.84$\,dB at $2.400$\,GHz after a $4.9$\,dB loss from L2 to \PL. The same
freedom optimized on L1 loses $8.4$\,dB to \PL{} ($-14.12$ to $-5.69$\,dB).
These designs also detune by about $75$\,MHz between models, more than the
copper designs, because the thickness change adds a detune of its own.
The copper-only design meets the target at \PL{} with less freedom; on
this evidence that is due to the frequency correction removing the
transfer loss, not to copper being a better design variable than
thickness.

\paragraph{What the data do not separate.}
We have not compared the adjoint-driven optimizer with a gradient-free
search or with a low-dimensional parametric sweep of patch size and feed
position on the same models. The final design widens the patch and the
feed, which a parametric sweep could also explore, so the result does not
show that the $38$-parameter basis or the shape gradient is necessary for
this antenna.

\section{Limitations and Future Work}
\label{sec:limits}

\paragraph{Feed.}
The feed is a uniform lumped-port strip across the $1.6$\,mm gap. The
final design's $4.3$\,mm port is a strip of that width, not a coaxial pin
(an SMA centre pin is $1.27$\,mm). A real feed would need a wide post or
strip, or a matching structure presenting the same port, and none is
modelled or verified. The copper outline is manufacturable on a standard
board; the feed is not shown to be.

\paragraph{Convergence and centring.}
\PL{} is the finest model and it is not converged. Centring the design on
the extrapolated converged response would take another space-mapping step
with a larger shift, but the $12.7$--$48.4$\,MHz spread of the
extrapolations exceeds the design's $11$\,MHz of total slack, so the
extrapolation would have to be narrowed first, for example with a
second-order L2 solve. The adjoint is first-order only; a $p = 2$ shape
adjoint would let the optimizer run on the more accurate model directly.

\paragraph{Scope.}
The objective is impedance match only; radiation pattern, gain and
efficiency are not computed. Conductors are PEC, so copper loss is not
modelled. The design is one antenna on one radius and has not been
fabricated or measured.

\paragraph{Comparison with structured-grid methods.}
A natural question is how a structured-grid design method would fare on
the same curved board at matched accuracy. A Meep \citep{oskooi2010meep}
model of this fixture has been set up and piloted, but its runs are not
yet converged, so we make no cost or accuracy comparison. Our working
hypothesis is that a body-fitted mesh reaches a given accuracy on the
curved conductor at lower cost than a staircased grid; testing it needs a
converged FDTD reference, ideally including a conformal FDTD scheme
\citep{dey1997locally,zagorodnov2003uniformly}, and is left to future work.

\section{How geode-fem Was Designed and Built}
\label{sec:built}

geode-fem implements the whiteroom L1--L4 operator specification of
Crutcher Dunnavant \citep{dunnavant2026whiteroom}. He also gave design
input on the solver's operator layer and is the author of 36 commits, made
in June 2026, that migrated the code to Rust 2024, added the Metal (Apple
GPU) backend, restructured the core crate and the examples into their
present layout, and created the geode-util crate. The code, the
experiments and the drafts of this paper were otherwise produced by
frontier AI models working as role-separated agents: Anthropic's Claude
(Opus, Fable and Sonnet models, through Claude Code) and OpenAI's GPT Sol
and GPT Astra models (via Codex). Codex reviewed the figures and page
layout of a draft of this paper. The agents were coordinated by Loom
\citep{loom}, an issue, worktree and pull-request workflow, and by a shared
Squad room.

The five pull requests that produced this paper's results (the fixture,
the smooth basis, the refinement study, the thickness ablation and the
copper design) each received an independent review by a separate agent
session, which recomputed reported numbers from the committed artifacts
with its own scripts, reproduced forward points and mutation-tested the
regression pins. Several claims were corrected because of those reviews.
Three examples follow. The converged resonance of the initial patch was
first bracketed $29$--$43$\,MHz below the band from a single fitted slope;
solving each level triple for its own order showed the order falling, and
the result is now a lower bound with an undetermined upper end
(Section~\ref{sec:convergence}). A review found that re-basing the feed
between two regenerated meshes changes the copper by $0.09$ or $0.47$\,mm,
so the sibling-mesh check and the L2-chain results are reported as
perturbed or related designs (Section~\ref{sec:results}). And a review
required the paper to state which way the converged answer probably falls,
rather than calling it unknown.

Counts from repository history, up to this revision: the paper and its
benchmark artifacts changed through 13 merged pull requests from
2026-07-20 and 30 commits. Of these commits, 26 carry a Claude co-author
trailer (Opus 19, Fable 7); the other 4 carry no trailer and contain about
$9\%$ of the added text lines. The human author directed the research,
made every scoping and claim decision, and takes responsibility for the
content.

\section{Reproducibility}
\label{sec:repro}

Every result in this paper is read from a committed, scripted
artifact.\footnote{The geode-fem source is public at
\url{https://github.com/rjwalters/geode-fem}. Archival DOI to be added at
deposit.}
% TODO(operator): Zenodo DOI (deposit stays an unpublished draft until approved)
Table~\ref{tab:artifacts} lists them, relative to
\path{benchmarks/patch_antenna_conformal/}. Each TOML records the exact
command, the mesh SHA-256 and the commit that produced it, and regression
tests pin the reported values. The meshes are regenerated from YAML
recipes in \path{reference/gmsh/} with gmsh. The figure scripts in
\path{figures/src/} read the artifacts directly, and \path{provenance.md}
in this paper's directory maps every number to its source field.

\begin{table}[t]
  \centering
  \caption{Artifacts behind each result, relative to
  \texttt{benchmarks/patch\_antenna\_conformal/}.}
  \label{tab:artifacts}
  \footnotesize
  \begin{tabular}{@{}>{\raggedright\arraybackslash}p{0.30\linewidth}>{\raggedright\arraybackslash}p{0.64\linewidth}@{}}
    \toprule
    Result & Artifact \\
    \midrule
    Fixture and flat-versus-bent baseline & \path{baseline_2g4_L0.toml}, \path{baseline_2g4_L1.toml} \\
    Refinement study, absorber ladder, cost (Section~\ref{sec:convergence}, Figure~\ref{fig:convergence}) & \path{convergence.toml}, \path{convergence/}, \path{upml_ladder_L1.toml} \\
    Copper optimizations (Table~\ref{tab:chain}) & \path{copper/*.toml}, \path{copper/README.md} \\
    Verifications and bands (Table~\ref{tab:verify}, Figure~\ref{fig:response}) & \path{copper/verify/}, \path{copper/g2_summary.toml} from \path{scripts/copper_g2.py} \\
    Frequency shift fit & \path{scripts/copper_shift_fit.py} \\
    Geometry of the designs (Figure~\ref{fig:geometry}) & \path{geometry/rbD_L1_copper_38_sm45.json} and the other \path{geometry/*copper*.json} \\
    Thickness ablation (Section~\ref{sec:ablations}) & \path{nested_dof/}, \path{nested_dof/verify/}, \path{scripts/nested_dof_g2.py} \\
    Generator and tests & \path{crates/geode-core/examples/patch_conformal_nested_dof.rs}; \path{crates/geode-core/tests/patch_conformal_copper.rs} \\
    \bottomrule
  \end{tabular}
\end{table}

\section{Conclusion}
\label{sec:conclusion}

A differentiable finite-element model of a curved, PML-terminated patch
antenna, with a feed that moves, gives exact shape gradients at one
factorization per frequency, and with them a copper-only design on a
standard $1.6$\,mm board reaches $-10.78$\,dB across the ISM band on the
finest model we can compute. The harder part was the discretization. The
resonance is not converged on any affordable mesh, an L1 optimum
lands $45$--$48$\,MHz off at \PL, and the margin left at the
lower band edge ($6.2$\,MHz) is smaller than every estimate of the
remaining shift. A frequency-only space mapping removed most of the
transfer loss, and the same data say the converged design probably misses
at $2.400$\,GHz. For a gradient-designed antenna, the refinement study is
part of the result.

\paragraph{Acknowledgments.}
We thank Crutcher Dunnavant for the whiteroom specification, for his
design input to the solver's operator layer and for his contributions to
the codebase. The AI-orchestrated development of geode-fem and of this
paper is described in Section~\ref{sec:built}.
% TODO(operator): approve acknowledgment wording

\bibliographystyle{plainnat}
\bibliography{refs}

\end{document}
